\section{Metodología}
\label{sec:metodologia}

Esta sección documenta \emph{qué} resuelve el programa y \emph{cómo}. Todas las
ecuaciones que siguen son las que están implementadas, no una descripción
genérica: donde el código se aparta de la forma de libro se dice y se explica por
qué.

\subsection{Qué se resuelve, y sobre qué dominio}

El programa integra en el tiempo cinco bloques acoplados sobre la misma malla:

\begin{enumerate}[leftmargin=*,itemsep=2pt]
\item \textbf{Momentum} del gas, con el lecho tratado como medio poroso.
\item \textbf{Energía}, una sola temperatura para gas y sólido (equilibrio
térmico local).
\item \textbf{Especies gaseosas}: CO, CO$_2$, H$_2$, H$_2$O, CH$_4$, N$_2$ y
O$_2$.
\item \textbf{Química heterogénea}, que consume e imprime en trece fases
sólidas.
\item \textbf{Submodelos diagnósticos}: cohesión, hinchamiento y magnetismo, que
leen los campos pero no realimentan al solucionador.
\end{enumerate}

\noindent
El dominio es el crisol entero ---metal, cavidad, tapa y carga---, no sólo el
lecho. Cada celda lleva una etiqueta y las propiedades varían con ella: no hay
un solucionador para el lecho y otro para el gas, hay uno con coeficientes que
cambian en el espacio.

\subsection{¿Todo esto se resuelve de verdad en 3-D?}
\label{sec:es3d}

Sí, y conviene ser explícito porque es una pregunta legítima: un problema con
simetría de revolución podría resolverse en 2-D axisimétrico o incluso en 1-D, y
entonces llamarlo «3-D» sería una etiqueta.

\begin{clave}
Los campos son arreglos de $n_x \times n_y \times n_z$ celdas y \textbf{las tres
componentes del gradiente, de la divergencia y del laplaciano se evalúan en las
tres direcciones}, con flujos en las seis caras de cada celda. No se impone
simetría axial, no se resuelve un sector, no hay ninguna dirección promediada.
La malla del preajuste grueso es $16\times16\times35 = \num{8960}$ celdas; la
media, $31\times31\times70 = \num{67270}$; la fina,
$62\times62\times139 = \num{534316}$.
\end{clave}

\noindent
De esas \num{8960} celdas, 308 son lecho, \num{3559} gas de la cavidad,
\num{1027} pared del crisol, 342 tapa y \num{3724} exterior. El solucionador
trabaja sobre las activas ---lecho y gas--- y las de metal entran sólo en la
conducción de calor.

\noindent
Ahora bien, «resuelto en 3-D» y «tridimensionalmente relevante» son cosas
distintas, y la honestidad exige separarlas. Cada término de las ecuaciones se
calcula en 3-D, pero no todos \emph{importan} en este ensayo. El programa mide
la magnitud de cada uno en cada paso y la reporta:

\begin{table}[H]
\centering\small
\begin{tabular}{lll}
\toprule
\textbf{Término} & \textbf{Máximo medido en la corrida} & \textbf{¿Domina?} \\
\midrule
Difusión de calor & --- & \textbf{sí, siempre} \\
Advección de calor & $\mathrm{Pe}_T^{\max} = \num{3.93}$ & sólo en el pico de gas \\
Difusión de especies & --- & \textbf{sí, siempre} \\
Advección de especies & $\mathrm{Pe}_m^{\max} = \num{5.88}$ & sólo en el pico de gas \\
Arrastre de Darcy & $\tau_D = \SI{1.8}{\micro\second}$ & \textbf{sí, en el lecho} \\
Viscoso de Brinkman & --- & sí, en el gas libre \\
Advección de momentum & $\mathrm{Re}_{\text{celda}}^{\max} = \num{4.51}$ & no \\
Forchheimer & $\mathrm{Re}_p^{\max} = \num{0.394}$ & \textbf{no} \\
Boyancia & $\mathrm{Ra}^{\max} = \num{1862}$ vs $\mathrm{Ra}_c = \num{1708}$ & \ojo{véase abajo} \\
Reacción heterogénea & --- & \textbf{sí} \\
Gasificación del char & $\mathrm{Da} = \num{1.6e-17}$ & \textbf{no} \\
\bottomrule
\end{tabular}
\caption{Todos los términos se evalúan en tres dimensiones, con flujos en las
seis caras de cada celda. Los máximos son sobre las 145 instantáneas de la
corrida, no valores nominales.}
\label{tab:terminos3d}
\end{table}

\ojo{Una advertencia honesta sobre la boyancia.} El caso la desactiva
(\texttt{activar\_boyancia: false}) con la justificación de que el Rayleigh del
ensayo está por debajo del umbral convectivo. La corrida \emph{contradice} esa
justificación: $\mathrm{Ra}$ llega a \num{1862} en el pico de devolatilización,
por encima del crítico \num{1708}. El término está implementado y basta cambiar
una línea del caso para activarlo; que no se haya hecho es una decisión
pendiente, no un resultado. Dicho eso, el pico es breve y el flujo en ese
momento lo domina la fuente de gas, tres órdenes por encima de cualquier
corriente convectiva, así que no se espera que cambie las conclusiones ---pero
eso hay que comprobarlo, no afirmarlo.

\paragraph{Advección y difusión son comparables, no despreciables.} Los Péclet
llegan a unidades en el estallido de gas: durante esos segundos el gas generado
arrastra calor y especies más rápido de lo que difunden. Una versión anterior de
este trabajo afirmaba lo contrario y hubo que corregirla; la prueba automática
que lo fija comprueba ahora la forma real ---un pico breve y valores
despreciables en el estado final--- y no una desigualdad global.

\paragraph{Qué aporta el 3-D que un 0-D no puede.} Ésta es la justificación real
de resolverlo así, y es verificable en los resultados:

\begin{itemize}[leftmargin=*,itemsep=2pt]
\item \textbf{Los gradientes térmicos dentro del lecho.} El 0-D tiene una sola
temperatura; aquí la pared, el fondo y el techo del lecho van por caminos
distintos, y de eso depende dónde y cuándo cuaja el aglomerado.
\item \textbf{El crecimiento del aglomerado.} Con temperatura uniforme todas las
celdas cruzarían el umbral de coquización a la vez y el cuerpo aparecería de
golpe. Con el campo 3-D real crece.
\item \textbf{La geometría del crisol.} El collar, la tapa y las paredes
inclinadas fijan el camino térmico y el volumen libre; nada de eso existe en un
modelo agregado.
\item \textbf{El campo de composición.} Las fases no son uniformes en el lecho,
y el visor muestra esa distribución.
\end{itemize}

\subsection{Momentum: Darcy--Brinkman--Forchheimer}

Una sola ecuación cubre el lecho granular y el espacio libre:

\begin{equation}
\frac{\rho}{\varepsilon}\frac{\partial \mathbf{u}}{\partial t}
+ \frac{\rho}{\varepsilon^{2}}(\mathbf{u}\cdot\nabla)\mathbf{u}
= -\nabla P
+ \mu_{\mathrm{ef}}\nabla^{2}\mathbf{u}
\underbrace{- \frac{\mu}{K}\mathbf{u}}_{\text{Darcy}}
\underbrace{- \frac{\rho C_F}{\sqrt{K}}|\mathbf{u}|\mathbf{u}}_{\text{Forchheimer}}
+ \rho\,\mathbf{g}\,\beta\,(T-T_{\mathrm{ref}}),
\label{eq:momentum}
\end{equation}

con $\mathbf{u} = (u,v,w)$ y todos los operadores en tres dimensiones. En el gas
libre $\varepsilon = 1$ y $K \rightarrow \infty$: los dos términos porosos
desaparecen y \eqref{eq:momentum} se reduce a Navier--Stokes puro.

El acoplamiento presión--velocidad es la \textbf{proyección de Chorin} sobre
malla escalonada (MAC), en tres pasos:

\begin{align}
\frac{\mathbf{u}^{*}-\mathbf{u}^{n}}{\Delta t}
&= -\frac{1}{\varepsilon}(\mathbf{u}^{n}\!\cdot\!\nabla)\mathbf{u}^{n}
+ \frac{\varepsilon}{\rho}\Big[\mu_{\mathrm{ef}}\nabla^{2}\mathbf{u}
- \frac{\rho C_F}{\sqrt K}|\mathbf{u}|\mathbf{u}
+ \rho\mathbf{g}\beta\,\Delta T\Big]
- \frac{\mathbf{u}}{\tau_D},
\label{eq:predictor}\\[4pt]
\nabla^{2}\phi &= \frac{\rho}{\Delta t}\,\nabla\!\cdot\!\mathbf{u}^{*} - \frac{\rho}{\Delta t} S_m,
\qquad \left.\frac{\partial \phi}{\partial n}\right|_{\partial\Omega} = 0,
\label{eq:poisson}\\[4pt]
\mathbf{u}^{n+1} &= \mathbf{u}^{*} - \frac{\Delta t}{\rho}\nabla\phi,
\qquad P^{n+1} = P^{n} + \phi .
\label{eq:correccion}
\end{align}

Aquí $\tau_D = \rho K/(\mu\varepsilon)$ es el tiempo de relajación de Darcy y
$S_m$ la fuente volumétrica de masa por devolatilización. La malla escalonada
---$u$ en caras $x$ de forma $(n_x{+}1, n_y, n_z)$, $v$ en caras $y$, $w$ en
caras $z$, escalares en centros--- evita el modo de presión en tablero de
ajedrez.

\paragraph{Compatibilidad del problema de Neumann.} Con \eqref{eq:poisson}
puramente Neumann, la solución existe sólo si el término fuente tiene media
nula. Con fuente de masa neta no la tiene: \emph{está demostrado con una prueba
que un recinto cerrado con fuente neta no admite solución}. Por eso el venteo no
es un adorno sino una condición de existencia. La incompatibilidad residual se
mide y se reporta ($\leq\SI{2.2e-16}{\per\second}$).

\subsection{Energía}

Una sola temperatura, con la capacidad térmica efectiva de cada celda:

\begin{equation}
(\rho c_p)_{\mathrm{ef}}\frac{\partial T}{\partial t}
= -\,(\rho c_p)_{\mathrm{gas}}\;\mathbf{u}\cdot\nabla T
+ \nabla\!\cdot\!\big(k_{\mathrm{ef}}\nabla T\big)
+ \dot{q}_{\mathrm{rxn}} + \dot{q}_{\mathrm{pir}},
\label{eq:energia}
\end{equation}

donde $(\rho c_p)_{\mathrm{ef}} = (1-\varepsilon)(\rho c_p)_{s} + \varepsilon(\rho c_p)_{g}$.

Dos detalles de \eqref{eq:energia} no son cosméticos y costaron una corrección
(\S\ref{sec:defectos}):

\begin{itemize}[leftmargin=*,itemsep=2pt]
\item \textbf{El término advectivo va en forma advectiva}, no en divergencia.
Como
$\nabla\!\cdot\!(\mathbf{u}T) = \mathbf{u}\cdot\nabla T + T\,\nabla\!\cdot\!\mathbf{u}$
y aquí $\nabla\!\cdot\!\mathbf{u}\neq 0$ ---la devolatilización crea unos
\qty{19}{\centi\meter\cubed\per\second} de gas dentro del lecho---, el segundo
sumando no transporta nada: es dilatación. Medido, restaba \SI{-778}{\kelvin\per\second}
de media al lecho. El gas que aparece sale del sólido \emph{a la temperatura
local} y no puede enfriar la celda que lo genera. En el código se calcula la
divergencia y se le restituye $T\,\nabla\!\cdot\!\mathbf{u}$, lo que deja
exactamente $\mathbf{u}\cdot\nabla T$; con campo solenoidal ambas formas
coinciden, y por eso las pruebas MMS no lo detectaban.
\item \textbf{Advecta el gas, no el bulto.} El coeficiente es
$(\rho c_p)_{\mathrm{gas}}/(\rho c_p)_{\mathrm{ef}}$, que vale 1 en el gas libre
y $\sim\!\num{2e-3}$ en el lecho. Es la formulación estándar de equilibrio
térmico local en medio poroso; usar la capacidad efectiva equivaldría a que el
gas arrastrase la entalpía del sólido, 500 veces mayor de la que puede llevar.
\end{itemize}

\subsection{Especies}

Para cada especie gaseosa, en concentración molar por volumen de celda:

\begin{equation}
\frac{\partial c_i}{\partial t}
= -\nabla\!\cdot\!(\mathbf{u}\,c_i)
+ \nabla\!\cdot\!\big(D_{i,\mathrm{ef}}\nabla c_i\big)
+ \dot{\omega}_i ,
\qquad i \in \{\mathrm{CO},\mathrm{CO_2},\mathrm{H_2},\mathrm{H_2O},\mathrm{CH_4},\mathrm{N_2},\mathrm{O_2}\},
\end{equation}

con difusividades de Fuller--Schettler--Giddings y corrección de tortuosidad
dentro del lecho:

\begin{equation}
D_i(T,P) = D_{i,\mathrm{ref}}\left(\frac{T}{T_{\mathrm{ref}}}\right)^{1.75}
\frac{P_{\mathrm{ref}}}{P},
\qquad
D_{i,\mathrm{ef}} = \varepsilon^{1.5}\, D_i ,
\qquad
D_{i,\mathrm{ef}}\big|_{\text{metal}} = 0 .
\end{equation}

\subsection{Química}

La química no se reimplementa: \texttt{fisica/adaptador\_v3.py} es un puente a
los módulos ya validados del modelo 0-D, y la consistencia entre ambos se
verifica a \num{2e-16}.

\subsubsection{Reacciones}

\begin{table}[H]
\centering\small
\begin{tabular}{ll}
\toprule
\textbf{Bloque} & \textbf{Reacción} \\
\midrule
Secado & H$_2$O(l) $\rightarrow$ H$_2$O(g) \\
Devolatilización & volátil $\rightarrow$ CO, CO$_2$, H$_2$, H$_2$O, CH$_4$ \\
Reducción por CO & 3\,Fe$_2$O$_3$ + CO $\rightarrow$ 2\,Fe$_3$O$_4$ + CO$_2$ \\
 & Fe$_3$O$_4$ + CO $\rightarrow$ 3\,FeO + CO$_2$ \\
 & FeO + CO $\rightarrow$ Fe + CO$_2$ \\
 & FeTiO$_3$ + CO $\rightarrow$ Fe + TiO$_2$ + CO$_2$ \\
Reducción por H$_2$ & las mismas cuatro, con H$_2$/H$_2$O \\
Gasificación & C + CO$_2$ $\rightarrow$ 2\,CO \quad (Boudouard) \\
 & C + H$_2$O $\rightarrow$ CO + H$_2$ \quad (gas de agua) \\
Homogénea & CO + H$_2$O $\rightleftharpoons$ CO$_2$ + H$_2$ \quad (WGS) \\
Combustión & C + O$_2$ $\rightarrow$ CO/CO$_2$, con el reparto de Arthur \\
Escoria & 2\,FeO + SiO$_2$ $\rightarrow$ Fe$_2$SiO$_4$ \\
\bottomrule
\end{tabular}
\end{table}

\subsubsection{La ley de velocidad, y por qué no lleva tope de conversión}

Cada reducción se escribe como una constante de Arrhenius multiplicada por la
\emph{afinidad} termodinámica del gas:

\begin{equation}
r_j = \underbrace{A_j\,e^{-E_{a,j}/RT}}_{\text{cinética}}\;
\underbrace{n_{\text{óxido}}}_{\text{inventario}}\;
\underbrace{\max\!\left(0,\; 1 - \frac{Q}{K_{\mathrm{eq},j}(T)}\right)}_{\text{afinidad}}\;
\underbrace{f_{\mathrm{acc}}}_{\text{accesibilidad}},
\qquad
Q_{\mathrm{CO}} = \frac{p_{\mathrm{CO_2}}}{p_{\mathrm{CO}}},
\quad
Q_{\mathrm{H_2}} = \frac{p_{\mathrm{H_2O}}}{p_{\mathrm{H_2}}} .
\label{eq:afinidad}
\end{equation}

Ésta es la pieza central del modelo. Al reducir un óxido el gas se oxida, $Q$ se
acerca a $K_{\mathrm{eq}}$, la afinidad cae \emph{continuamente} a cero y la
reacción se para sola. \textbf{La conversión parcial emerge de la termodinámica
del gas}, sin ningún $\alpha_{\max}$ impuesto ---que es lo que hacía la versión
anterior del trabajo. De ahí sale el tamponamiento sobre la frontera
Fe$_3$O$_4$/FeO. Los dos canales, CO y H$_2$, se evalúan por separado y actúan a
la vez.

\begin{table}[H]
\centering\small
\begin{tabular}{lrrl}
\toprule
\textbf{Reacción} & $A$ & $E_a$ [\si{\kilo\joule\per\mole}] & \textbf{Fuente} \\
\midrule
Fe$_2$O$_3\rightarrow$Fe$_3$O$_4$ & \num{5.0e4} & 70{,}0 & Hammam \emph{et al.} (2022) \\
Fe$_3$O$_4\rightarrow$FeO & \num{2.0e5} & 88{,}7 & Hammam \emph{et al.} (2022), $E_{a1}$ \\
FeO$\rightarrow$Fe & \num{8.0e5} & 120{,}0 & Hammam (2022) y lit.\ de wüstita \\
FeTiO$_3\rightarrow$Fe+TiO$_2$ & \num{1.0e6} & 190{,}0 & extrapolado de 1000--\SI{1200}{\celsius} \\
Devolatilización & \num{1.0e4} & 55{,}0 & van Krevelen; compuerta a \SI{450}{\celsius} \\
Boudouard & \num{3.0e8} & 220{,}0 & lit.\ de cinética intrínseca de \emph{chars} \\
Gas de agua & \num{1.0e8} & 210{,}0 & ídem \\
WGS & \num{1.0e5} & 80{,}0 & forma reversible \\
Combustión & \num{1.0e6} & 150{,}0 & correlación de Arthur para CO/CO$_2$ \\
\bottomrule
\end{tabular}
\caption{Parámetros cinéticos. Los de reducción llevan además un factor de
escala por par redox declarado \textsc{calibrable}; el de la magnetita,
$k = 0{,}20$, queda pendiente de revisión (\S\ref{sec:frontera}).}
\end{table}

\paragraph{Compuerta de devolatilización.} La liberación de volátiles se activa
con una sigmoide en temperatura,
$g_{\mathrm{vol}} = \sigma(T - 273{,}15;\, T_0,\, w)$ con $w = \SI{40}{\celsius}$
y $T_0 = \SI{450}{\celsius}$. $T_0$ es el \emph{centro} de la sigmoide, no un
umbral: con el valor anterior de \SI{200}{\celsius} la compuerta estaba abierta
al \qty{98}{\percent} ya a \SI{360}{\celsius} (\S\ref{sec:puerta}).

\paragraph{Accesibilidad de las lamas.} La titanohematita está intercrecida
dentro de los granos de magnetita, así que sólo se expone al gas conforme avanza
el frente de reducción de la magnetita:
$f_{\mathrm{acc}} = f(\alpha_{\mathrm{Fe_3O_4}},\, e_{\mathrm{lama}})$, con el
espesor de lama declarado \textsc{calibrable} porque no se ha medido.

\subsubsection{Integración: por qué Patankar}

La química es rígida ---modos desde \SI{1e-9}{\second}--- y debe conservar la
positividad de las concentraciones: una concentración negativa no es un error
pequeño, es una imposibilidad física que además envenena el paso siguiente. Se
integra con un esquema de \textbf{Patankar}, que separa producción y consumo,

\begin{equation}
c_i^{\,n+1} = \frac{c_i^{\,n} + \Delta t\, P_i}{1 + \Delta t\, D_i / c_i^{\,n}},
\qquad
\frac{\mathrm{d}c_i}{\mathrm{d}t} = P_i - D_i ,
\end{equation}

y es \textbf{incondicionalmente positivo}, en subpasos internos acotados por
\texttt{dt\_quimica\_max\_s} (\SI{0.05}{\second}) con un máximo de 256.

\paragraph{Tabulación.} La termodinámica depende sólo de $T$, así que se tabula
una vez sobre una rejilla y se interpola. En malla fina eso multiplicó por
\textbf{202} el número de celdas resueltas por segundo.

\subsection{Cierres}

\begin{table}[H]
\centering\small
\begin{tabular}{lll}
\toprule
\textbf{Magnitud} & \textbf{Relación} & \textbf{Fuente} \\
\midrule
Permeabilidad & $K = \dfrac{d_p^{2}\,\varepsilon^{3}}{150\,(1-\varepsilon)^{2}}$ & Kozeny--Carman \\[7pt]
Conductividad del lecho & $k_{\mathrm{ef}} = (1-\varepsilon)k_{\mathrm{esq}} + \varepsilon k_{\mathrm{gas}}$ & mezcla volumétrica \\[3pt]
Capacidad del lecho & $(\rho c_p)_{\mathrm{ef}} = (1-\varepsilon)(\rho c_p)_s + \varepsilon(\rho c_p)_g$ & --- \\[3pt]
Densidad del gas & $\rho = PM/RT$ & gas ideal \\[3pt]
Difusividades & $D \propto T^{1.75}/P$, $D_{\mathrm{ef}} = \varepsilon^{1.5}D$ & Fuller \emph{et al.} (1966) \\[3pt]
Termoquímica & Shomate y Maier--Kelley & NIST--JANAF (Chase, 1998) \\
\bottomrule
\end{tabular}
\end{table}

\noindent
Para el concentrado del ensayo ($d_p = \SI{175}{\micro\meter}$,
$\varepsilon = 0{,}54$) Kozeny--Carman da $K = \SI{1.52e-10}{\meter\squared}$.
La conductividad del lecho interpola entre \qty{1.2}{\watt\per\meter\per\kelvin}
y \qty{0.6}{} según el grado de reducción local (Kiamehr \emph{et al.}).

\paragraph{La porosidad está viva.} $\varepsilon$ se recalcula en cada paso a
partir del volumen de sólido que queda,

\begin{equation}
\varepsilon(t) = 1 - \big(1-\varepsilon_0\big)\,
\frac{\sum_i c_i(t)\,M_i/\rho_i}{\sum_i c_i(0)\,M_i/\rho_i},
\end{equation}

y con ella se refrescan $K$, $k_{\mathrm{ef}}$ y la tortuosidad de las
difusividades. Se aplica el cambio \emph{relativo} sobre la porosidad inicial
calibrada, de modo que en $t=0$ sale exactamente el valor del ensayo. Sólo se
actualiza cuando ha cambiado más de \num{5e-3}: reescribirla en cada paso
invalidaba la caché de la factorización ILU y hacía la corrida diez veces más
lenta.

\subsection{Condiciones de frontera}

\begin{itemize}[leftmargin=*,itemsep=3pt]
\item \textbf{Mufla.} Radiación sobre la superficie exterior del crisol,
\begin{equation}
q'' = \varepsilon_{\mathrm{em}}\,F\,\sigma\,\big(T_{\mathrm{mufla}}^{4}(t) - T^{4}\big),
\qquad \varepsilon_{\mathrm{em}} = 0{,}80,\quad F = 1{,}00,
\end{equation}
con la \emph{curva real medida} del horno, incluida su caída al abrir la puerta.
El factor de vista es 1,00: el acoplamiento térmico no lleva ningún parámetro
calibrado.
\item \textbf{Venteo.} La devolatilización genera unos
\qty{19}{\centi\meter\cubed\per\second} de gas dentro de un recinto que, si se
tratara como cerrado, no tendría solución (véase la nota de
\eqref{eq:poisson}). Se modela como un sumidero de masa en la corona de gas bajo
la junta de la tapa, y el inventario que sale se contabiliza para cerrar el
balance.
\item El resto de las caras son de flujo difusivo nulo, $\partial\phi/\partial n = 0$.
\end{itemize}

\subsection{Submodelos diagnósticos}

Tres, y los tres \emph{leen} los campos sin realimentar al solucionador:

\begin{itemize}[leftmargin=*,itemsep=2pt]
\item \textbf{Cohesión}: coquización del carbón a través de una variable interna
irreversible, el grado de plastificación, que avanza mientras la celda está
dentro de la ventana termoplástica (350--\SI{500}{\celsius}, van Krevelen) y no
vuelve atrás. El aglomerado es el conjunto de celdas conexas con
$c \geq 0{,}5$, obtenido por componentes conexas en 3-D.
\item \textbf{Hinchamiento} (\S\ref{sec:hinchamiento}): sigue al mismo grado de
plastificación, no a la temperatura instantánea.
\item \textbf{Magnetismo} (\S\ref{sec:magnetismo}): magnetización de saturación
a temperatura ambiente del inventario de fases, mezcla lineal en masa.
\end{itemize}

\subsection{Método numérico}

\paragraph{Separación de operadores de Strang.} Los procesos difieren en órdenes
de magnitud ---relajación de Darcy \SI{1.8e-6}{\second}, química con modos desde
\SI{1e-9}{\second}, advección $\sim\SI{1e-2}{\second}$---, así que se integran
por bloques:

\begin{equation}
\text{química}\!\left(\tfrac{\Delta t}{2}\right) \rightarrow
\text{momentum}(\Delta t) \rightarrow
\text{transporte}(\Delta t) \rightarrow
\text{química}\!\left(\tfrac{\Delta t}{2}\right).
\end{equation}

Es simétrico y de segundo orden \emph{para la separación}, pero conviene no
engañarse: el orden global lo limita el bloque menos preciso, y la verificación
por MMS mide \num{1.01} en el tiempo. La simetría sigue mereciendo la pena
porque elimina el sesgo sistemático de la separación de Lie.

\paragraph{Volúmenes finitos conservativos.} Los flujos se evalúan en caras, de
modo que lo que sale de una celda entra exactamente en su vecina. La difusión
promedia \textbf{conductividades} con media armónica,

\begin{equation}
k_{i+1/2} = \frac{2\,k_i\,k_{i+1}}{k_i + k_{i+1}},
\end{equation}

y cada fila se divide por la capacidad de su propia celda. Promediar
difusividades es correcto sólo con $\rho c_p$ uniforme; en la interfaz gas/metal
el salto es de 3.100 veces y esa aproximación creaba energía
(\S\ref{sec:defectos}).

\paragraph{Advección TVD.} El esquema por omisión es \textbf{TVD superbee} con
el limitador de Roe--Sweby,

\begin{equation}
\psi(r) = \max\!\big[0,\ \min(2r,1),\ \min(r,2)\big],
\qquad
r = \frac{\phi_{U} - \phi_{UU}}{\phi_{D} - \phi_{U}},
\end{equation}

que degrada localmente a upwind junto a las fronteras, donde no hay segunda
celda aguas arriba. Se comparó con upwind y con diferencias centradas:

\begin{table}[H]
\centering\small
\begin{tabular}{lccc}
\toprule
$\mathrm{Pe}$ & \textbf{upwind} & \textbf{central} & \textbf{TVD superbee} \\
\midrule
1 & $L_2 = \num{2.2e-3}$ & \num{3.1e-4} & \num{3.0e-4} \\
10 & \num{3.1e-2} & \num{5.1e-3} & \num{5.1e-3} \\
50 & rango $[0;\,0{,}444]$ & \ojo{$[-0{,}049;\,0{,}444]$} & $[0;\,0{,}444]$ \\
\bottomrule
\end{tabular}
\caption{A $\mathrm{Pe}=50$ el esquema central produce concentraciones
\emph{negativas}. TVD mantiene monotonía con la misma precisión.}
\end{table}

\paragraph{Qué se trata implícitamente, y por qué.} El arrastre de Darcy es
lineal en la velocidad, así que se integra de forma \emph{exacta} e
incondicionalmente estable:

\begin{equation}
u^{n+1} = \frac{u^{n} + \Delta t\,(\text{resto})}{1 + \Delta t/\tau_D} .
\end{equation}

\begin{table}[H]
\centering\small
\begin{tabular}{lrr}
\toprule
\textbf{Configuración} & \textbf{paso estable} & \textbf{pasos para \SI{720}{\second}} \\
\midrule
Todo explícito & \SI{3.92e-7}{\second} & \num{1836734694} \\
+ Darcy implícito & \SI{5.43e-5}{\second} & \num{13248000} \\
+ viscoso implícito & \SI{5.43e-5}{\second} & \num{13248000} \\
+ difusión implícita & \SI{2.15e-3}{\second} & \num{335520} \\
\bottomrule
\end{tabular}
\caption{Ganancia total: \num{5474} veces. Sin Darcy implícito,
$\tau_D = \SI{1.8}{\micro\second}$ obligaría a mil ochocientos millones de pasos.}
\end{table}

\paragraph{Sistemas lineales.} BiCGSTAB con precondicionador ILU, con la
factorización cacheada por $(\Delta t, \nu, \text{Darcy})$, y respaldo directo
(\texttt{spsolve}) por si el iterativo sufre \emph{breakdown} ---que ocurre
cuando la física llega al reposo y $\|b\|\rightarrow 0$, es decir, justamente
cuando el resultado es correcto. Los sistemas se construyen \textbf{sólo sobre
las celdas activas}: meter las celdas sólidas, que llevan
$\varepsilon = 10^{-6}$ y $K = 10^{-20}$, ponía catorce órdenes de contraste en
la misma matriz. Restringirlo llevó la divergencia residual de \num{1.2e-3} a
\num{1.7e-17} y además resultó un \qty{48}{\percent} más rápido.

\paragraph{Paso adaptativo.} Se limita por CFL ($\mathrm{CFL} = 0{,}4$), por
estabilidad difusiva y por el acoplamiento químico, con
$\Delta t \in [\SI{1e-9}{}, \SI{0.5}{}]\,\si{\second}$. En régimen el limitador
práctico es \texttt{dt\_max\_acople}.

\subsection{Geometría y malla}

El crisol se reconstruyó de las cotas y se cerró con una fotografía acotada, que
aportó una cota ausente de todos los archivos: un collar de
\qty{30.6}{\milli\meter} de diámetro, más ancho que la boca. Con él la masa
calculada coincide con los \qty{32.67}{\gram} declarados a
\SI{4e-14}{\gram}; sin él faltaban \qty{77}{\milli\gram}.

El lecho es un disco delgado ---\qty{1.36}{\centi\meter\cubed} en
\qty{22.8}{\milli\meter} de diámetro son \qty{3.26}{\milli\meter} de altura, una
relación de aspecto de 7,7:1---, y por eso la malla es \textbf{anisótropa}:
$\Delta x = \Delta y = \SI{2.0}{\milli\meter}$ y $\Delta z =
\SI{1.0}{\milli\meter}$ en el preajuste grueso, la mitad en el medio y la cuarta
parte en el fino. Las celdas cortadas por la frontera llevan fracción
volumétrica parcial, de modo que la carga se reparte con su volumen real y no
por el centro de celda.

\subsection{Qué produce}

Cada instantánea es un NPZ con el contrato declarado en
\texttt{docs/CONTRATOS.md}: velocidades en caras, presión, temperatura,
porosidad, cohesión e hinchamiento en centros, las siete especies gaseosas, las
trece fases sólidas, las etiquetas de celda y unos metadatos con los números
adimensionales y el magnetismo. De ahí salen el visor 3-D y \textbf{todas} las
cifras de este informe.
